\chapter{Testy i weryfikacja}
Rozdział przedstawia sposób sprawdzenia poprawności działania aplikacji oraz ocenę uzyskanych rezultatów. Testy skupiają się na weryfikacji poprawności działania oraz zgodności z wymaganiami funkcjonalnymi i niefunkcjonalnymi.

\section{Testy funkcjonalne}
\begin{table}[ht!]
\centering
\caption{Tabela wyników testów funkcjonalnych}
\label{tab:wyniki}
\begin{tabular}{lll}
\toprule
Test                            & Oczekiwany Wyniki             & Wynik \\  
\midrule
Połączenie z API                & Dane zwracane w JSON          & Pozytywny \\
Wyświetlanie aktualnej temp.    & Widoczna poprawna wartość     & Pozytywny \\ 
Wyświetlanie dziennika          & Lista rekordów                & Pozytywny \\
\bottomrule 
\end{tabular}
\sourcetab{opracowanie własne}
\end{table}

\section{Testy niefunkcjonalne}
Test wydajności
\begin{itemize}
	\item Czas odpowiedzi API: akceptowalny
	\item Aplikacja reaguje płynnie. \\
\end{itemize}

Test kompatybilności
\begin{itemize}
	\item Android 8, 10, 13 – działanie poprawne.
\end{itemize}

\section{Wyniki testów}
System działa poprawnie i spełnia wszystkie założenia funkcjonalne. Komunikacja między ESP32, bazą danych i aplikacją mobilną przebiega bez błędów. Dane są poprawnie zapisywane i odczytywane. Interfejs użytkownika jest czytelny i intuicyjny.
